\section{Conclusion}
\label{sec:conclusion}

We presented \StereoSplat, a feed-forward framework that reconstructs
metric-scale dynamic 3D Gaussian scenes from calibrated stereo video
streams. By anchoring scale to the stereo baseline and focal length
through the triangulation formula $Z = b f / d$, and by replacing the
orthographic canonical space of \StreamSplat with a perspective
rasterizer and a metric back-projection head, we obtain a representation
whose Gaussians live in the metric world frame. A streaming display
module with motion-adaptive EMA and opacity-stability filtering addresses
the temporal flickering and popping artifacts that plague dynamic GS
methods. Experiments on TartanAir, DSEC, and KITTI Stereo confirm both
metric scale recovery and competitive photometric quality, and a
five-axis ablation isolates the contribution of each component.

\paragraph{Limitations and Future Work.}
The metric accuracy of \StereoSplat is bounded by the underlying stereo
matcher; in low-texture and occluded regions the disparity is unreliable
and the metric depth degrades, although our confidence masking mitigates
the worst failures. The current pipeline assumes a calibrated rectified
stereo pair; extending to uncalibrated or non-rig stereo (e.g., two
independently moving cameras) is an open direction. We currently rely on
an external SLAM trajectory for moving-rig poses; jointly estimating
poses and dynamic Gaussians in a single feed-forward pass is a natural
next step, and would connect this work to the event-driven setting of
EventGS where scale recovery from a single event camera is even more
challenging.
